% Algebra 2 Lesson 6-2 fast-track packet: Exponential Models (enVision A2, Example 2 and Example 4 only)
% Kept for Savvas Lesson Quiz items 3, 4, 5 (compound interest, continuous compounding, comparing options).
% Cut: Example 1 (monthly rate of a population; quiz item 1), Example 3 (the natural base e; its one fact is
% stated before Example 4), Example 5 (model from two points; quiz item 2), Example 6 (regression; no quiz item).
% Shared preamble for the "fast" (hacked-down) Algebra 2 lesson packets, Oct 2026.
% Each packet: % Shared preamble for the "fast" (hacked-down) Algebra 2 lesson packets, Oct 2026.
% Each packet: % Shared preamble for the "fast" (hacked-down) Algebra 2 lesson packets, Oct 2026.
% Each packet: \input{fast-packet-template} then \begin{document} ... \end{document}.
\documentclass[12pt]{article}
\usepackage[margin=0.7in]{geometry}
\usepackage{amsmath,amssymb}
\usepackage{booktabs,tabularx,array}
\usepackage{enumitem}
\usepackage[T1]{fontenc}
\usepackage[utf8]{inputenc}
\usepackage{lmodern}
\usepackage{tikz}
\usepackage{pgfplots}
\pgfplotsset{compat=1.18}
\usepackage{xcolor}
\setlist{nosep,leftmargin=1.4em}
\setlength{\parindent}{0pt}
\setlength{\parskip}{6pt}
\renewcommand{\arraystretch}{1.35}
\newcolumntype{Y}{>{\raggedright\arraybackslash}X}
\pagestyle{empty}

% blanks and workspace
\newcommand{\blank}[1][2.6cm]{\rule[-2pt]{#1}{0.5pt}}
\newcommand{\lines}[1][1]{\par\vspace{2pt}\foreach \i in {1,...,#1}{\rule{0pt}{22pt}\rule{\linewidth}{0.4pt}\par}}
\newcommand{\wspart}[2]{\vspace{4pt}{\large\bfseries #1: #2}\par\vspace{-6pt}\rule{\textwidth}{0.4pt}\par}

% header: \packethead{Lesson 1-3}{Piecewise-Defined Functions}
\newcommand{\packethead}[2]{%
{\LARGE\bfseries #2}\hfill Name \blank[5cm]\\[2pt]
Algebra 2 \quad #1 (fast track) \hfill Date \blank[2.5cm] \quad Period \blank[1.2cm]\par\vspace{2pt}}

% boxed goals: \goals{learning goal}{language goal}
\newcommand{\goals}[2]{%
\fbox{\parbox{\dimexpr\textwidth-2\fboxsep-2\fboxrule}{%
\textbf{Learning goal:} #1\\[2pt]
\textbf{Language goal:} #2}}\par}

% quiz target box: \quiztarget{text}
\newcommand{\quiztarget}[1]{%
\vspace{4pt}\colorbox{gray!12}{\parbox{\dimexpr\textwidth-2\fboxsep}{\textbf{On the quiz:} #1}}\par}

% IXL table row helper: code & skill & which quiz question it feeds
\newcommand{\ixlrow}[3]{\texttt{#1} & #2 & #3 \\}

% one grid style for every graph; override xmin/xmax/... per plot
\pgfplotsset{sheetgrid/.style={
  axis lines=middle, grid=both, grid style={gray!45},
  xmin=-10.5, xmax=10.5, ymin=-10.5, ymax=10.5,
  xtick={-10,-8,...,10}, ytick={-10,-8,...,10},
  x=0.28cm, y=0.28cm,
  ticklabel style={font=\scriptsize, fill=white, inner sep=1pt},
  xlabel={$x$}, ylabel={$y$},
  axis line style={thick}, clip=true}}
% unit grid for small windows
\pgfplotsset{unitgrid/.style={sheetgrid,
  xmin=-5.5, xmax=5.5, ymin=-5.5, ymax=5.5,
  xtick={-5,...,5}, ytick={-5,...,5}, x=0.5cm, y=0.5cm}}
% open and closed endpoints
\tikzset{closed/.style={circle, fill, inner sep=1.6pt},
         open/.style={circle, draw, fill=white, inner sep=1.4pt}}
 then \begin{document} ... \end{document}.
\documentclass[12pt]{article}
\usepackage[margin=0.7in]{geometry}
\usepackage{amsmath,amssymb}
\usepackage{booktabs,tabularx,array}
\usepackage{enumitem}
\usepackage[T1]{fontenc}
\usepackage[utf8]{inputenc}
\usepackage{lmodern}
\usepackage{tikz}
\usepackage{pgfplots}
\pgfplotsset{compat=1.18}
\usepackage{xcolor}
\setlist{nosep,leftmargin=1.4em}
\setlength{\parindent}{0pt}
\setlength{\parskip}{6pt}
\renewcommand{\arraystretch}{1.35}
\newcolumntype{Y}{>{\raggedright\arraybackslash}X}
\pagestyle{empty}

% blanks and workspace
\newcommand{\blank}[1][2.6cm]{\rule[-2pt]{#1}{0.5pt}}
\newcommand{\lines}[1][1]{\par\vspace{2pt}\foreach \i in {1,...,#1}{\rule{0pt}{22pt}\rule{\linewidth}{0.4pt}\par}}
\newcommand{\wspart}[2]{\vspace{4pt}{\large\bfseries #1: #2}\par\vspace{-6pt}\rule{\textwidth}{0.4pt}\par}

% header: \packethead{Lesson 1-3}{Piecewise-Defined Functions}
\newcommand{\packethead}[2]{%
{\LARGE\bfseries #2}\hfill Name \blank[5cm]\\[2pt]
Algebra 2 \quad #1 (fast track) \hfill Date \blank[2.5cm] \quad Period \blank[1.2cm]\par\vspace{2pt}}

% boxed goals: \goals{learning goal}{language goal}
\newcommand{\goals}[2]{%
\fbox{\parbox{\dimexpr\textwidth-2\fboxsep-2\fboxrule}{%
\textbf{Learning goal:} #1\\[2pt]
\textbf{Language goal:} #2}}\par}

% quiz target box: \quiztarget{text}
\newcommand{\quiztarget}[1]{%
\vspace{4pt}\colorbox{gray!12}{\parbox{\dimexpr\textwidth-2\fboxsep}{\textbf{On the quiz:} #1}}\par}

% IXL table row helper: code & skill & which quiz question it feeds
\newcommand{\ixlrow}[3]{\texttt{#1} & #2 & #3 \\}

% one grid style for every graph; override xmin/xmax/... per plot
\pgfplotsset{sheetgrid/.style={
  axis lines=middle, grid=both, grid style={gray!45},
  xmin=-10.5, xmax=10.5, ymin=-10.5, ymax=10.5,
  xtick={-10,-8,...,10}, ytick={-10,-8,...,10},
  x=0.28cm, y=0.28cm,
  ticklabel style={font=\scriptsize, fill=white, inner sep=1pt},
  xlabel={$x$}, ylabel={$y$},
  axis line style={thick}, clip=true}}
% unit grid for small windows
\pgfplotsset{unitgrid/.style={sheetgrid,
  xmin=-5.5, xmax=5.5, ymin=-5.5, ymax=5.5,
  xtick={-5,...,5}, ytick={-5,...,5}, x=0.5cm, y=0.5cm}}
% open and closed endpoints
\tikzset{closed/.style={circle, fill, inner sep=1.6pt},
         open/.style={circle, draw, fill=white, inner sep=1.4pt}}
 then \begin{document} ... \end{document}.
\documentclass[12pt]{article}
\usepackage[margin=0.7in]{geometry}
\usepackage{amsmath,amssymb}
\usepackage{booktabs,tabularx,array}
\usepackage{enumitem}
\usepackage[T1]{fontenc}
\usepackage[utf8]{inputenc}
\usepackage{lmodern}
\usepackage{tikz}
\usepackage{pgfplots}
\pgfplotsset{compat=1.18}
\usepackage{xcolor}
\setlist{nosep,leftmargin=1.4em}
\setlength{\parindent}{0pt}
\setlength{\parskip}{6pt}
\renewcommand{\arraystretch}{1.35}
\newcolumntype{Y}{>{\raggedright\arraybackslash}X}
\pagestyle{empty}

% blanks and workspace
\newcommand{\blank}[1][2.6cm]{\rule[-2pt]{#1}{0.5pt}}
\newcommand{\lines}[1][1]{\par\vspace{2pt}\foreach \i in {1,...,#1}{\rule{0pt}{22pt}\rule{\linewidth}{0.4pt}\par}}
\newcommand{\wspart}[2]{\vspace{4pt}{\large\bfseries #1: #2}\par\vspace{-6pt}\rule{\textwidth}{0.4pt}\par}

% header: \packethead{Lesson 1-3}{Piecewise-Defined Functions}
\newcommand{\packethead}[2]{%
{\LARGE\bfseries #2}\hfill Name \blank[5cm]\\[2pt]
Algebra 2 \quad #1 (fast track) \hfill Date \blank[2.5cm] \quad Period \blank[1.2cm]\par\vspace{2pt}}

% boxed goals: \goals{learning goal}{language goal}
\newcommand{\goals}[2]{%
\fbox{\parbox{\dimexpr\textwidth-2\fboxsep-2\fboxrule}{%
\textbf{Learning goal:} #1\\[2pt]
\textbf{Language goal:} #2}}\par}

% quiz target box: \quiztarget{text}
\newcommand{\quiztarget}[1]{%
\vspace{4pt}\colorbox{gray!12}{\parbox{\dimexpr\textwidth-2\fboxsep}{\textbf{On the quiz:} #1}}\par}

% IXL table row helper: code & skill & which quiz question it feeds
\newcommand{\ixlrow}[3]{\texttt{#1} & #2 & #3 \\}

% one grid style for every graph; override xmin/xmax/... per plot
\pgfplotsset{sheetgrid/.style={
  axis lines=middle, grid=both, grid style={gray!45},
  xmin=-10.5, xmax=10.5, ymin=-10.5, ymax=10.5,
  xtick={-10,-8,...,10}, ytick={-10,-8,...,10},
  x=0.28cm, y=0.28cm,
  ticklabel style={font=\scriptsize, fill=white, inner sep=1pt},
  xlabel={$x$}, ylabel={$y$},
  axis line style={thick}, clip=true}}
% unit grid for small windows
\pgfplotsset{unitgrid/.style={sheetgrid,
  xmin=-5.5, xmax=5.5, ymin=-5.5, ymax=5.5,
  xtick={-5,...,5}, ytick={-5,...,5}, x=0.5cm, y=0.5cm}}
% open and closed endpoints
\tikzset{closed/.style={circle, fill, inner sep=1.6pt},
         open/.style={circle, draw, fill=white, inner sep=1.4pt}}

\begin{document}

\packethead{Lesson 6-2}{Exponential Models}
\goals{I can write an exponential model for money that earns compound interest or continuously compounded interest, and use it to find the value of the account.}{I can explain why more compounding periods give a larger account value, using the words principal, interest rate, compounding periods, compounded continuously, and natural base e.}

\quiztarget{(Savvas numbers) (1) A town's yearly decrease as a rate per quarter and (2) a point on $y=a\cdot b^x$ from two points: full Savvas quiz only. You will (3) build $A(t)=P\left(1+\frac rn\right)^{nt}$ for an account compounded monthly (Example 2); (4) find the value of an account compounded continuously (Example 4); (5) bonus: compare simple, continuous and yearly interest on \$8{,}000 and say which option is linear and which are exponential (Examples 2 and 4, Practice).}

\wspart{Example 2}{Understand Compound Interest}
\textbf{Compound interest} is interest earned on interest already deposited. \textbf{Compound interest formula:} $A=P\left(1+\frac rn\right)^{nt}$, where $P=$ the initial \textbf{principal} invested, $r=$ annual \textbf{interest rate} as a decimal, $n=$ number of \textbf{compounding periods} per year, and $A=$ the value of the account after $t$ years.

Tamira invests \$5,000 in an account that pays 4\% annual interest. \textbf{A.} How much will there be in the account after 3 years if the interest is compounded annually, quarterly, or monthly? ($P=5000$, $r=0.04$, $t=3$.)

{\centering\small\renewcommand{\arraystretch}{1.12}
\begin{tabular}{|l|c|l|r|}
\hline
 & $n$ & \textbf{Compound interest formula} & \textbf{After 3 years} \\ \hline
Annually & 1 & $A=5000\left(1+\frac{0.04}{1}\right)^{3(1)}=5000(1.04)^{3}$ & \$5,624.32 \\
Quarterly & 4 & $A=5000\left(1+\frac{0.04}{4}\right)^{3(4)}=5000(1.01)^{12}$ & \$5,634.13 \\
Monthly & 12 & $A=5000\left(1+\frac{0.04}{12}\right)^{3(12)}\approx5000(1.003333)^{36}$ & \$5,636.36 \\ \hline
\end{tabular}\par}
As the number of compounding periods increases, the amount in the account also increases: the more often interest is added, the earlier it earns interest of its own.

% TE transcription error corrected: a2_6-2_TE.tex line 256 prints the monthly base as (1+0.003)^{12}; the TE's own note
% (line 258) says 0.003333..., which gives the printed 1.0407. The packet uses 1.003333 in the monthly rows of both tables.
\textbf{B.} What is the annual growth rate for each possible account? Use the Power of a Power rule, $\left(b^{n}\right)^{t}=b^{nt}$, to rewrite each model with exponent $t$. Growth factor minus 1 is the rate.

{\centering\small\renewcommand{\arraystretch}{1.12}
\begin{tabular}{|l|l|c|}
\hline
 & \textbf{Rewrite with exponent} $t$ & \textbf{Annual growth rate} \\ \hline
Annually & $A=5000\left(1+\frac{0.04}{1}\right)^{1t}=5000(1.04)^{t}$ & 4\% \\
Quarterly & $A=5000\left(1+\frac{0.04}{4}\right)^{4t}=5000\left((1.01)^{4}\right)^{t}\approx5000(1.0406)^{t}$ & 4.06\% \\
Monthly & $A=5000\left(1+\frac{0.04}{12}\right)^{12t}\approx5000\left((1.003333)^{12}\right)^{t}\approx5000(1.0407)^{t}$ & 4.07\% \\ \hline
\end{tabular}\par}
Quarterly: $1.0406-1=0.0406=4.06\%$. This yearly rate is the \textbf{effective annual interest rate}.

\textbf{Say it:} Compounded monthly, the account has \blank[1.2cm] compounding periods per year instead of 1, so its value after 3 years is \blank[2cm] than with annual compounding, because \dots\lines[1]

\newpage

\wspart{Try It 2}{Compound interest}
\$3,000 is invested in an account that earns 3\% annual interest, compounded monthly.

$P=\blank[1.6cm]$ \quad $r=\blank[1.6cm]$ \quad $n=\blank[1.2cm]$ \qquad $A=\blank[1.6cm]\left(1+\dfrac{\blank[1.4cm]}{\blank[1cm]}\right)^{\blank[1.2cm]\,t}$

\textbf{a.} What is the value of the account after 10 years? 100 years?

\begin{minipage}[t]{0.48\textwidth}
After 10 years:\lines[3]
\end{minipage}\hfill
\begin{minipage}[t]{0.48\textwidth}
After 100 years:\lines[3]
\end{minipage}

\textbf{b.} What is the effective annual interest rate of the account?

Hint: rewrite the model with exponent $t$, as in Part B.\lines[2]

\wspart{Example 4}{Find Continuously Compounded Interest}
\textbf{The natural base $e$.} Put \$1 in at 100\% interest for one year, compounded $n$ times: the value is $\left(1+\frac1n\right)^n$. For $n=1$ it is 2; $n=100$, 2.704813829; $n=100{,}000$, 2.718268237. As $n$ keeps growing, the value approaches the irrational number $e=2.718281828459\ldots$, the \textbf{natural base}. Compounding that never stops is \textbf{compounded continuously}.

\textbf{Continuously compounded interest formula:} $A=Pe^{rt}$, where $P=$ the principal, $e=$ the natural base, $r=$ annual interest rate as a decimal, and $A=$ the value after $t$ years.

\begin{minipage}[t]{0.56\textwidth}
\vspace{0pt}
Regina invests \$12,600 in an account that is compounded continuously, at a 3.2\% interest rate for a 12-year term. What is the value of the account after 12 years? Round your answer to the nearest dollar.
\begin{enumerate}[label=\textbf{\arabic*.}, itemsep=2pt]
\item $P=12{,}600$, \ $r=0.032$, \ $t=12$.
\item $A=Pe^{rt}=12{,}600\,e^{0.032(12)}$
\item Multiply the exponent first: $0.032(12)=0.384$, so $A=12{,}600\,e^{0.384}$.
\item Use the $e^x$ key: $e^{0.384}\approx1.468145$, so $A\approx18{,}498.63$.
\end{enumerate}
\textbf{To the nearest dollar, the value is \$18,499.}
\end{minipage}\hfill
\begin{minipage}[t]{0.40\textwidth}
\vspace{0pt}
\fbox{\parbox{\dimexpr\linewidth-2\fboxsep-2\fboxrule}{\small\raggedright\textbf{Calculator trap.} When you evaluate $e^{0.032(12)}$, either simplify $0.032(12)$ as $0.384$ first, or type parentheses, \texttt{e\^{}(0.032*12)}, so that $e$ is raised to the entire product, not just the first factor.}}
\end{minipage}

\textbf{Say it:} The account is \blank[3.2cm], so I use $A=\blank[1.6cm]$. The base $e$ is the \blank[2.6cm], about \blank[1.6cm].

\newpage

\wspart{Try It 4}{Continuously compounded interest}
You invest \$125,000 in an account that earns 4.75\% annual interest, compounded continuously.

$P=\blank[2cm]$ \quad $r=\blank[1.6cm]$ \qquad $A=\blank[2cm]\,e^{\blank[1.6cm]\,t}$

\textbf{a.} What is the value of the account after 15 years?\\
\textbf{b.} What is the value of the account after 30 years?

\begin{minipage}[t]{0.48\textwidth}
\textbf{a.} After 15 years:\lines[3]
\end{minipage}\hfill
\begin{minipage}[t]{0.48\textwidth}
\textbf{b.} After 30 years:\lines[3]
\end{minipage}

Why did you use $A=Pe^{rt}$ and not $A=P\left(1+\frac rn\right)^{nt}$? Use the words \emph{compounded continuously} and \emph{natural base}.\lines[2]

% R5 Practice: structural parallel of Savvas Lesson Quiz item 5 (Paul, $8,000, 6 years, 5% simple /
% 4.5% continuous / 4.75% yearly). Same structure, new numbers and setting. Key: A 6600, B 6615.65, C 6712.35;
% B - A = $16, C - B = $97 (96.70 to the nearest dollar); year 1: A 5200, B 5178.10, C 5187.50 (linear gains most at first).
\wspart{Practice}{Three ways to save}
Lena has \$5,000 from a summer job to save for 8 years. She compares three accounts.

\begin{minipage}[t]{0.58\textwidth}
\vspace{0pt}
\textbf{Option A:} 4\% simple interest\\
\textbf{Option B:} 3.5\% interest compounded continuously\\
\textbf{Option C:} 3.75\% interest compounded yearly
\end{minipage}\hfill
\begin{minipage}[t]{0.39\textwidth}
\vspace{0pt}
\small\textbf{Simple interest} is paid on the principal only: $A=P(1+rt)$. It adds the same amount, $Pr$, every year.
\end{minipage}

\textbf{a.} Write a model for each option and fill in the table. Round to the cent.

\begin{center}
\begin{tabular}{|l|l|c|c|}
\hline
 & \textbf{Model} & \textbf{After 1 year} & \textbf{After 8 years} \\ \hline
A (simple) & $A=\rule{0pt}{18pt}\blank[3.6cm]$ & \blank[2.2cm] & \blank[2.2cm] \\
B (continuous) & $A=\rule{0pt}{18pt}\blank[3.6cm]$ & \blank[2.2cm] & \blank[2.2cm] \\
C (yearly) & $A=\rule{0pt}{18pt}\blank[3.6cm]$ & \blank[2.2cm] & \blank[2.2cm] \\ \hline
\end{tabular}
\end{center}

\textbf{b.} After 8 years, to the nearest dollar, Option B will be worth \$\blank[1.4cm] more than Option A and \$\blank[1.4cm] less than Option C.

\textbf{c.} Option A is a \blank[2.4cm] function, while Options B and C are \blank[2.4cm] functions. Which option gained the most in the first year? Which gains faster over time? Use your table.\lines[2]

\newpage
\wspart{IXL}{Practice for the quiz}
Go to IXL, Algebra 2. Type the code in the search box. Do the starred skills first.

\vspace{2pt}
\begin{tabularx}{\textwidth}{l Y Y}
\toprule
\textbf{Code} & \textbf{Skill} & \textbf{Feeds} \\
\midrule
\ixlrow{LWE}{Evaluate exponential functions (warm-up)}{Do Now: calculator powers and the e key}
\ixlrow{YJW $\bigstar$}{Compound interest: word problems}{Quiz 3 and 5 (Example 2)}
\ixlrow{5GU $\bigstar$}{Continuously compounded interest: word problems}{Quiz 4 and 5 (Example 4)}
\ixlrow{KLF $\bigstar$}{Describe linear and exponential growth and decay}{Quiz 5 (bonus), last parts: linear or exponential}
\bottomrule
\end{tabularx}

\vspace{2pt}
\textbf{SmartScore 70 or above on each starred ($\bigstar$) skill replaces your quiz score.}

\wspart{Exit question}{Who will have more?}
Adam and Jacinta open an account on the same day. Each starts with \$8,000 for 20 years. Adam's interest is compounded quarterly at 1.25\%. Jacinta's interest is compounded continuously at 1.25\%. Who do you predict will have more money in their account after 20 years? Explain your reasoning.

\vspace{4pt}
Adam: $A=\blank[5.5cm]$ \qquad Jacinta: $A=\blank[5cm]$

\vspace{6pt}
\blank[2.2cm] will have more, because interest that is \blank[3.6cm] is added more often than interest with only \blank[1.2cm] compounding periods per year, so \dots\lines[3]

\vspace{6pt}
\textbf{Before the quiz, finish these sentences.}

The \blank[2.6cm] is the money you put in at the start; $r$ is the annual \blank[3cm], written as a decimal.

More \blank[4cm] per year means a \blank[2cm] account value.

When interest is \blank[4cm], I use $A=Pe^{rt}$, where $e\approx\blank[1.6cm]$ is the \blank[3cm].

\end{document}
